\begin{figure}[htbp]
\centering
\begin{tikzpicture}[font=\small,>=Stealth]
\draw[fill=gray!6] (0,0) circle(1.4);
\draw[dashed,gray] (-1.4,0) arc[start angle=180,end angle=0,x radius=1.4,y radius=.4];
\draw[gray] (-1.4,0) arc[start angle=180,end angle=360,x radius=1.4,y radius=.4];
\draw[thick,<->] (-1.7,-1.275)--(1.7,1.275);
\fill (0,0) circle(.04) node[below right] {$0$};
\foreach \x/\y/\lab/\pos in {1.12/.84/x/above left,-1.12/-.84/{-x}/below right}{
\fill[blue!75!black] (\x,\y) circle(.075);\node[\pos] at (\x,\y) {$\lab$};}
\node[left] at (-1.4,.65) {$S^2$};
\node[above right] at (1.7,1.275) {line $\R x$};
\draw[->,thick] (2.7,0)--(4.2,0) node[midway,above] {$\pi_S$};
\fill[blue!75!black] (4.8,0) circle(.075);
\node[right,align=left] at (5,0) {one point $[x]=[-x]$\\of $\mathbb{RP}^2$};
\node[align=center] at (2.5,-2) {a line through the origin $=$ an antipodal pair $=$ a projective point};
\end{tikzpicture}
\caption{The sphere model of projective space pairs $x$ only with $-x$.
The matching filled markers lie on the same line through the origin.
The final dot represents one class, not a drawing of the whole quotient.}
\label{fig:projective-sphere}
\end{figure}
